% =========================================================
\section{Contexto y problema}
% =========================================================

\begin{frame}{Contexto general}
  \begin{columns}[T,onlytextwidth]
    \begin{column}{0.50\textwidth}
      \begin{UCBlock}{Idea central}
        Inserte aquí el contexto general del área de investigación, la relevancia académica o aplicada del tema y el escenario donde se sitúa la tesis de posgrado.
      \end{UCBlock}
      \begin{itemize}
        \item Antecedente principal del problema.
        \item Brecha o necesidad que motiva el estudio.
        \item Alcance general del proyecto de Magíster/Doctorado.
      \end{itemize}
    \end{column}
    \begin{column}{0.46\textwidth}
      \centering
      \begin{tikzpicture}[scale=0.92,every node/.style={font=\small}]
        \fill[UCBlueLight] (0,0) rectangle (6.3,4.15);
        \draw[UCBluePPT,line width=0.8pt] (0,0) rectangle (6.3,4.15);
        \node[draw=UCBluePPT,fill=UCWhite,rounded corners=2pt,align=center,text width=1.85cm] (a) at (1.15,3.25) {Variable\\de entrada};
        \node[draw=UCBluePPT,fill=UCWhite,rounded corners=2pt,align=center,text width=1.72cm] (b) at (3.25,3.25) {Proceso\\o método};
        \node[draw=UCBluePPT,fill=UCWhite,rounded corners=2pt,align=center,text width=1.70cm] (c) at (5.25,3.25) {Condición\\de análisis};
        \node[draw=UCYellowDeep,fill=UCYellow!18!white,rounded corners=2pt,align=center,text width=4.25cm] (d) at (3.15,1.65) {Sistema de estudio\\o caso de aplicación};
        \node[draw=UCBlueDark,fill=UCBlueDark,text=UCWhite,rounded corners=2pt,align=center,text width=2.55cm] (e) at (3.15,0.42) {Resultado\\de ejemplo};
        \draw[-{Latex[length=2mm]},UCBluePPT,line width=0.8pt] (a) -- (d);
        \draw[-{Latex[length=2mm]},UCBluePPT,line width=0.8pt] (b) -- (d);
        \draw[-{Latex[length=2mm]},UCBluePPT,line width=0.8pt] (c) -- (d);
        \draw[-{Latex[length=2mm]},UCYellowDeep,line width=0.9pt] (d) -- (e);
      \end{tikzpicture}
      {\scriptsize Figura de ejemplo: reemplazar por esquema conceptual propio.}
    \end{column}
  \end{columns}
\end{frame}

\begin{frame}{Motivación}
  \begin{columns}[T,onlytextwidth]
    \begin{column}{0.48\textwidth}
      \begin{UCBlock}{Motivación académica}
        Inserte aquí por qué el tema es relevante para la literatura, la disciplina o la formación de conocimiento en ingeniería.
      \end{UCBlock}
      \begin{UCNeutralBox}
        Consejo: conecta la motivación con una brecha verificable, no solo con una descripción general del área.
      \end{UCNeutralBox}
    \end{column}
    \begin{column}{0.48\textwidth}
      \begin{UCBlock}{Motivación aplicada}
        Inserte aquí la relevancia práctica, industrial, social, tecnológica o metodológica del proyecto.
      \end{UCBlock}
      \UCsmallmetric{Impacto}{Variable de ejemplo}{Indicar el aporte esperado en una frase breve.}
      \UCsmallmetric{Escala}{Caso de ejemplo}{Indicar si el estudio es conceptual, experimental, computacional o aplicado.}
    \end{column}
  \end{columns}
\end{frame}

\begin{frame}{Problema de investigación}
  \begin{columns}[T,onlytextwidth]
    \begin{column}{0.52\textwidth}
      \begin{UCBlock}{Problema}
        Inserte aquí el problema de investigación. Debe quedar claro qué no se conoce, qué no está resuelto o qué requiere ser evaluado.
      \end{UCBlock}
      \begin{itemize}
        \item Elemento 1 del problema.
        \item Elemento 2 del problema.
        \item Elemento 3 del problema.
      \end{itemize}
    \end{column}
    \begin{column}{0.44\textwidth}
      \begin{UCKeyPoint}
        \textbf{Mensaje clave.} La lámina debe permitir que la comisión entienda por qué la tesis de posgrado es necesaria y qué decisión metodológica se deriva del problema.
      \end{UCKeyPoint}
    \end{column}
  \end{columns}
\end{frame}

\begin{frame}{Pregunta de investigación}
  \begin{columns}[T,onlytextwidth]
    \begin{column}{0.52\textwidth}
      \begin{UCBlock}{Pregunta principal}
        ¿Cómo influye \textit{variable de ejemplo} sobre \textit{resultado de ejemplo} en \textit{sistema o caso de estudio} bajo \textit{condiciones de análisis}?
      \end{UCBlock}
    \end{column}
    \begin{column}{0.44\textwidth}
      \begin{UCKeyPoint}
        \textbf{Hipótesis o supuesto de trabajo.} Inserte aquí la hipótesis, con lenguaje prudente y verificable.
      \end{UCKeyPoint}
    \end{column}
  \end{columns}
\end{frame}
